% !PRINTED 74
\weberpage{74}
%<<CONT
\gesperrt{der einer Primzahl $q$ entsprechende Index $\beta$ aller Zahlen $b$ durch $q$ oder durch $q - 1$ theilbar ist, je nachdem $q$ mehrmals oder nur einmal unter den Primfactoren von $m$ vorkommt.}

\begin{center}
\S.\,20.

\gesperrt{Die Kreistheilungsperioden.}
\end{center}

Als die einfachsten Functionen, durch die man die Kreistheilungskörper darzustellen versuchen kann, bieten sich die Kreistheilungsperioden dar, die eine unmittelbare Verallgemeinerung der im \S.\,167 des ersten Bandes betrachteten Gauss'schen Perioden sind. Wir verstehen darunter Folgendes.

Es bedeute $\mathfrak{A}$ einen Theiler der Gruppe $\mathfrak{N}$ vom Index $e$ und $a$ durchlaufe die Zahlen von $\mathfrak{A}$. Ist $r$ eine primitive $m^{\text{te}}$ Einheitswurzel, so heisst die Summe
\weq{1}{\eta = \overset{a}{\Sigma}\, r^a}
\gesperrt{eine zu der Gruppe $\mathfrak{A}$ gehörige Kreistheilungsperiode} vom Index $e$.

Machen wir in (1) eine der Substitutionen $(r, r^a)$, so bleibt $\eta$ ungeändert, wie unmittelbar aus der Gruppeneigenschaft der $a$ folgt.

Um $\mathfrak{N}$ in die zu $\mathfrak{A}$ gehörigen Nebengruppen zu zerlegen, müssen wir die Zahlen $1, n_1, n_2, \ldots n_{e-1}$ aus $\mathfrak{N}$ passend auswählen, dass man
\weq{2}{\mathfrak{N} = \mathfrak{A} + \mathfrak{A} n_1 + \mathfrak{A} n_2 + \cdots + \mathfrak{A} n_{e-1}}
erhält. Machen wir dann in $\eta$ die Substitutionen
\weq{3}{(r, r), (r, r^{n_1}), \ldots (r, r^{n_{e-1}}),}
so geht $\eta$ in die conjugirten Perioden
\weq{4}{\eta, \eta_1, \ldots \eta_{e-1}}
über, und wenn diese alle von einander verschieden sind, so gehört $\eta$ zur Gruppe $\mathfrak{A}$. Der zu $\mathfrak{A}$ gehörige Körper $\Omega(\eta)$ besteht aus allen rationalen Functionen von $\eta$, und die Zahlen $\eta_1 \ldots \eta_{e-1}$ gehören alle zu derselben Gruppe $\mathfrak{A}$.

Wenn aber die Grössen (4) nicht alle von einander verschieden sind, so gehört die Zahl $\eta$ nicht zu der Gruppe $\mathfrak{A}$, sondern zu einer umfassenderen Gruppe $\mathfrak{A}'$, von der $\mathfrak{A}$ ein Theiler
%CONT>>
